\section{Regular Languages and Finite-State Automota}
\subsection{Deterministic Finite-State Automata (DFA)}

Suppose we want to describe a language \(L\) over \(\Sigma\), using only a finite amount of information. We could extend \(\Sigma\) with other symbols to an alphabet \(\Pi\) and assign (non-uniquely) a particular a word \(w \in \Pi^{\ast}\) to \(L\). But we can only do this for countably many \(L\). But we want \(w\) to \term{describe} \(L\).

\begin{example}
    For example, if \(\Sigma = \Bqty{0, 1}\) and \(L\) is all words over \(\Sigma\) that start with a \(1\). We can use \(\Pi = \Bqty{0, 1, a, \ldots, z, \_}\) and using natural language, \(w = \texttt{all words over \(\Sigma\) that start with a \(1\)}\).
\end{example}

We will see a countable class of languages, where we can do this not via natural language, but symbolically We first give a computational method to determine these languages.

\begin{definition}[Deterministic Finite-State Automaton (DFA)]
    A \term{deterministic finite-state automaton (DFA)} is a quintuple \(D = \pqty{Q, \Sigma, \delta, q_0, F}\) consisting of:
    \begin{itemize}
        \item a finite non-empty set of \term{states} \(Q\),
        \item an alphabet \(\Sigma\),
        \item a \term{transition function} \(\functiondomain{\delta}{Q \times \Sigma}{Q}\),
        \item a \term{starting state} \(q_0 \in Q\),
        \item a set of \term{accept states} \(F \subseteq Q\).
    \end{itemize}
\end{definition}

\begin{remark}
    The input to a DFA is any word \(w \in \Sigma^{\ast}\). The DFA starts in the starting state \(q_0\), and reads the first letter \(x_1\). It then moves to \(\delta\pqty{q_0, x_1}\) and reads \(x_2\), hence moving to \(\delta\pqty{\delta\pqty{q_0, x_1}, x_2}\), etc.

    When all letters of \(w\) is read, if the DFA is in an accept state, then \(w\) is \term{accepted}, and otherwise \(w\) is \term{rejected}.
\end{remark}

\begin{example*}
    The empty string \(\emptyword\) is accepted, if and only if \(q_0 \in F\).
\end{example*}

\begin{definition}[Transition Diagram]
    A \term{transition diagram} of a DFA \(D\) is a directed, labelled graph such that
    \begin{enumerate}[label=(\arabic*)]
        \item the vertex set is \(Q\),
        \item for each \(\pqty{q, \sigma} \in Q \times \Sigma\), draw an directed edge from \(q\) to \(\delta\pqty{q, \sigma}\) labelled by \(\sigma\). (This means there are exactly \(\abs{\Sigma}\) out-edges from every vertex.)
        \item add an arrow from ``nowhere'' to \(q_0\) to indicate the starting state,
        \item draw double circles around the accept states in \(F\) (and single circles otherwise).
    \end{enumerate}

    If there are lots of edges from some vertex to the same vertex, then we can draw one edge with label which is a subset of \(\Sigma\).
\end{definition}

\begin{example*}
    Suppose we work in the natural language alphabet, and we want to accept strings that contain the substring \texttt{cat}. One might attempt to draw a transition diagram as in \autoref{fig:dfa-cat-incorrect}.

    This is, however, incorrect, since it does not accept the string \texttt{ccat}, or \texttt{cacat}. The correct transition diagram is given in \autoref{fig:dfa-cat-correct}.
\end{example*}

\begin{figure}[ht!]
    \centering
    \begin{tikzpicture}[->,>=Stealth,shorten >=1pt,auto,node distance=3cm,semithick,initial text=Start]
        \node[state,initial]   (q0)                    {\(q_0\)};
        \node[state]           (q1) [right of=q0]      {\(q_1\)};
        \node[state]           (q2) [right of=q1]      {\(q_2\)};
        \node[state,accepting] (q3) [right of=q2]      {\(q_3\)};

        \path (q0) edge                node {\texttt{c}} (q1)
                (q1) edge                node {\texttt{a}} (q2)
                (q2) edge                node {\texttt{t}} (q3)
                (q0) edge [loop below]   node {\(\Sigma\setminus\Bqty{\texttt{c}}\)} (q0)
                (q1) edge [bend left=30] node {\(\Sigma\setminus\Bqty{\texttt{a}}\)} (q0)
                (q2) edge [bend right=30] node[above] {\(\Sigma\setminus\Bqty{\texttt{t}}\)} (q0)
                (q3) edge [loop right]  node {\(\Sigma\)} (q3);
    \end{tikzpicture}
    \caption{Incorrect Attempt of DFA for strings containing \texttt{cat}}
    \label{fig:dfa-cat-incorrect}
\end{figure}

\begin{figure}[ht!]
    \centering
    \begin{tikzpicture}[->,>=Stealth,shorten >=1pt,auto,node distance=3cm,semithick,initial text=Start]
        \node[state,initial]   (q0)                    {\(q_0\)};
        \node[state]           (q1) [right of=q0]      {\(q_1\)};
        \node[state]           (q2) [right of=q1]      {\(q_2\)};
        \node[state,accepting] (q3) [right of=q2]      {\(q_3\)};

        \path (q0) edge                node {\texttt{c}} (q1)
                (q1) edge                node {\texttt{a}} (q2)
                (q2) edge                node {\texttt{t}} (q3)
                (q0) edge [loop below]   node {\(\Sigma\setminus\Bqty{\texttt{c}}\)} (q0)
                (q1) edge [bend left=30] node {\(\Sigma\setminus\Bqty{\texttt{a, c}}\)} (q0)
                (q2) edge [bend right=30] node[above] {\(\Sigma\setminus\Bqty{\texttt{c, t}}\)} (q0)
                (q3) edge [loop right]  node {\(\Sigma\)} (q3)
                (q1) edge [loop below]   node {\texttt{c}} (q1)
                (q2) edge [bend left=30]   node {\texttt{c}} (q1);
    \end{tikzpicture}
    \caption{Correct Attempt of DFA for strings containing \texttt{cat}}
    \label{fig:dfa-cat-correct}
\end{figure}

\begin{definition}[Transition Table]
    A \term{transition table} of a DFA \(D\) is a table with rows indexed by the states \(Q\) and columns indexed by the alphabet \(\Sigma\), such that the entry in row \(q\) and column \(\sigma\) is \(\delta\pqty{q, \sigma}\).
\end{definition}

\subsection{Regular Languages}

\begin{definition}[Extended Transition Function (DFA)]
    \label{def:extended-transition-function-dfa}%
    The \term{extended transition function} of a DFA \(D\) is given by
    \[
        \functiondomain{\hat{\delta}}{Q \times \Sigma^{\ast}}{Q}
    \]
    defined inductively by
    \begin{itemize}
        \item \(\hat{\delta} \pqty{q, \epsilon} = q\),
        \item \(\hat{\delta} \pqty{q, \sigma_1} = \delta(q, \sigma_1)\), and
        \item \(\hat{\delta} \pqty{q, \sigma_1 \ldots \sigma_k} = \delta \pqty{\hat{\delta} \pqty{q, \sigma_1 \ldots \sigma_{k - 1}}, \sigma_k}\).
    \end{itemize}
\end{definition}

For any \(w \in \Sigma^{\ast}\), we have \(\hat{\delta}\pqty{q_0, w}\) is the end state of an input \(w\).

\begin{lemma}[Transitivity of Extended Transition Function]
    For all \(1 \leq l \leq k\), we have
    \[
        \hat{\delta}\pqty{q, \sigma_1 \ldots \sigma_k} = \hat{\delta}\pqty{\hat{\delta}\pqty{q, \sigma_1 \ldots \sigma_l}, \sigma_{l + 1} \ldots \sigma_k}.
    \]
\end{lemma}

\begin{proof}
    We may show by induction on \(k\). Base case \(k = 1\) is trivial. Suppose the result holds for \(k - 1\). Then
    \begin{align*}
        \hat{\delta}\pqty{q, \sigma_1 \ldots \sigma_k} &= \delta\pqty{\hat{\delta}\pqty{q, \sigma_1 \ldots \sigma_{k - 1}}, \sigma_k} \\
        &= \delta\pqty{\hat{\delta}\pqty{\hat{\delta}\pqty{q, \sigma_1 \ldots \sigma_l}, \sigma_{l + 1} \ldots \sigma_{k - 1}}, \sigma_k} \\
        &= \hat{\delta}\pqty{\hat{\delta}\pqty{q, \sigma_1 \ldots \sigma_l}, \sigma_{l + 1} \ldots \sigma_k}.
    \end{align*}
\end{proof}

\begin{definition}[Regular Languages]
    Let \(D\) be a DFA. The \term{language accepted by \(D\)}, denoted by \(\langacc{L}(D)\), is the set of all words accepted by \(D\), i.e., the set of all \(w \in \Sigma^{\ast}\) such that \(\hat{\delta}\pqty{q_0, w} \in F\).
    
    A language \(L \subseteq \Sigma^{\ast}\) is \term{regular} if there exists a DFA \(D\) such that \(L = \langacc(D)\).
\end{definition}

\subsection{Non-Determinism and \(\emptyword\)-Transitions}

We now examine a (seemingly) more general construction, where we can choose which state to move to (or equivalently, moving to many states simultaneously).

\begin{definition}[Non-Deterministic Finite-State Automaton (NFA)]
    A \term{non-deterministic finite-state automaton (NFA)} is a quintuple \(N = \pqty{Q, \Sigma, \delta, q_0, F}\) consisting of:
    \begin{itemize}
        \item a finite non-empty set of \term{states} \(Q\),
        \item an alphabet \(\Sigma\),
        \item a \term{transition function} \(\functiondomain{\delta}{Q \times \Sigma}{\powerset\pqty{Q}}\),
        \item a \term{starting state} \(q_0 \in Q\),
        \item a set of \term{accept states} \(F \subseteq Q\).
    \end{itemize}
\end{definition}

\begin{remark}
    Note the difference of a NFA from a DFA is that the transition function \(\delta\) now maps to a \emph{set of states}, rather than a single state. This means that, for a given state and input symbol, the NFA can transition to a number of possible next states, possibly none at all (discussed below).
\end{remark}

An NFA has input \(w = \sigma_1, \ldots, \sigma_k \in \Sigma^{\ast}\) while in state \(q_0\), reads \(\sigma\), then ``moves to'' all states in \(\delta\pqty{q_0, \sigma_1} \subseteq Q\). For each of these states, it reads \(\sigma_2\) and moves to all states in \(\delta\pqty{q, \sigma_2}\), so we are now at
\[
    \bigcup_{p \in \delta\pqty{q_0, \sigma_1}} \delta\pqty{p, \sigma_2} \subseteq Q,
\]
etc. We continue to get ``end states'' \(E_W \subseteq Q\). If any of the resulting states after reading the entire input \(w\) is an accept state, i.e., \(E_w \cap F \neq \emptyset\), then \(w\) is accepted by the NFA, and otherwise \(w\) is rejected.

It could be well possible that \(\delta\pqty{p, \sigma} = \emptyset\) as \(\emptyset \in \powerset\pqty{Q}\). But what if we are currently in the set of states \(S \subseteq Q\), and we read \(\sigma\), only to find
\[
    \bigcup_{p \in S} \delta\pqty{p, \sigma} = \emptyset?
\]

Then we are in ``no state'' (the set of states is empty), and the NFA rejects the input \(w\).

Given an NFA \(N\), we may define the transition diagram and table in the same way as for a DFA.

\begin{figure}[ht!]
    \centering
    \begin{tikzpicture}[->,>=Stealth,shorten >=1pt,auto,node distance=3cm,semithick,initial text=Start]
        \node[state, initial] (q0)  {\(q_0\)};
        \node[state] (a)  [above right=of q0] {\(a\)};
        \node[state] (b)  [right=of q0] {\(b\)};
        \node[state] (c)  [right=of b] {\(c\)};
        \node[state] (d)  [below right=of q0] {\(d\)};

        \path (q0) edge [bend left=15] node {\(\sigma_1\)} (a)
            (q0) edge node {\(\sigma_1\)} (b)
            (q0) edge [bend left=20] node {\(\sigma_1\)} (d)
            (b) edge node {\(\sigma_2\)} (c)
            (d) edge node[swap] {\(\sigma_2\)} (c)
            (d) edge [bend left=20] node {\(\sigma_2\)} (q0)
            (d) edge [loop below] node {\(\sigma_2\)} (d);
    \end{tikzpicture}
    \caption{Part of an NFA}
    \label{fig:nfa-example}%
\end{figure}

\begin{example*}
    Consider \autoref{fig:nfa-example}. If we read the word \(w = \sigma_1 \sigma_2\), then after state one, we are in the set of states \(\Bqty{a, b, d}\). After reading \(\sigma_2\), we are in the set of states \(\Bqty{c, q_0, d}\).
\end{example*}

\begin{definition}[Extended Transtion Function (NFA)]
    \label{def:extended-transition-function-nfa}%
    The \term{extended transition function} of an NFA \(N\) is given by
    \[
        \functiondomain{\hat{\delta}}{Q \times \Sigma^{\ast}}{\powerset\pqty{Q}}
    \]
    defined inductively by
    \begin{itemize}
        \item \(\hat{\delta} \pqty{q, \epsilon} = \Bqty{q}\),
        \item \(\hat{\delta} \pqty{q, \sigma_1} = \delta(q, \sigma_1)\), and
        \item the inductive step
        \[
        \hat{\delta} \pqty{q, \sigma_1 \ldots \sigma_k} = \bigcup_{p \in \hat{\delta}\pqty{q, \sigma_1 \ldots \sigma_{k - 1}}} \delta\pqty{p, \sigma_k}.
        \]
    \end{itemize}
\end{definition}

In other words, \(\hat{\delta}\pqty{q, w}\) is the set of all ``end states'' \(E_w\)after reading \(w\) starting from \(q\).

Furthermore, it is going to be even more flexible, if we can change \emph{certain} states, without reading a symbol. Such transitions are called \term{\(\emptyword\)-transitions}. We would like to think of this as reading \(\emptyword\) (``nothing'') between the letters, whenever and as often as we like.

\begin{definition}[\(\emptyword\)-NFA]
    An \term{\(\emptyword\)-NFA} is a quintuple \(E = \pqty{Q, \Sigma, \delta, q_0, F}\) just like an NFA, except that the \term{transition function} is now
    \[
        \functiondomain{\delta}{Q \times \pqty{\Sigma \cup \Bqty{\emptyword}}}{\powerset\pqty{Q}}.
    \]
\end{definition}

Here, rather than the empty word being a symbol, it ``acts as a symbol''.

Now, we want to define how we process a word.

\begin{definition}[\(\emptyword\)-Closure]
    For an \(\emptyword\)-NFA \(E\) and a state \(Q\), we define the \term{\(\emptyword\)-closure} \(\ecl\pqty{q}\) of \(q\) as the set of states reached from \(q\), by repeated \(\emptyword\)-transitions. That is, define inductively \(S_0 \pqty{q} = \Bqty{q}\), \(S_1 \pqty{q} = \Bqty{q} \cup \delta\pqty{q, \emptyword}\), and 
    \[
        S_{i + 1} \pqty{q} = S_i \pqty{q} \cup \bigcup_{r \in S_i\pqty{q}} \delta\pqty{r, \emptyword}.
    \]
    and
    \[
        \ecl\pqty{q} = \bigcup_{i = 0}^{\infty} S_i\pqty{q} = S_N\pqty{q} \subseteq Q.
    \]
    for some \(N\) (since this is an ascending chain of finite sets, which are subsets of a finite set).

    The \term{\(\emptyword\)-closure} of a set of states \(S \subseteq Q\) is defined as
    \[
        \ecl\pqty{S} = \bigcup_{q \in S} \ecl\pqty{q}.
    \]

    Clearly we have \(S \subseteq \ecl\pqty{S}\), and the \(\emptyword\)-closure of a union of sets is the union of the \(\emptyword\)-closures of the sets.

    A given set of state \(R \subseteq Q\) is \term{\(\emptyword\)-closed} if \(\ecl\pqty{R} = R\).
\end{definition}

\begin{lemma}[\(\emptyword\)-Closure is \(\emptyword\)-Closed]
    For any \(S \subseteq Q\), \(\ecl\pqty{S}\) is \(\emptyword\)-closed.
\end{lemma}

\begin{proof}
    Let \(C = \ecl\pqty{S}\). We would like to show that
    \[
        \ecl\pqty{C} \subseteq C = \ecl\pqty{S}.
    \]

    Let \(x \in \ecl\pqty{C}\). Then \(x \in \ecl\pqty{r}\) for some \(r \in C\). So there exists some finite sequence of \(\emptyword\)-transitions from \(r\) to \(x\), i.e. mathematically, we have \(r_0 = r\), \(r_1 \in \delta\pqty{r_0, \emptyword}\), \(r_2 \in \delta\pqty{r_1, \emptyword}\), \dots, \(r_n = x \in \delta\pqty{r_{n - 1}, \emptyword}\).

    But \(r \in \ecl\pqty{S}\), so there is some \(s \in S\) with length \(m\) sequence of \(\emptyword\)-transitions from \(s\) to \(r\). We put these together to get a sequence of \(\emptyword\)-transitions from \(s\) to \(x\), which means \(x \in S_{m + n} \pqty{s} \subseteq \ecl\pqty{s} \subseteq \ecl\pqty{S}\).
\end{proof}

An \(\emptyword\)-NFA has input \(w = \sigma_1 \cdots \sigma_k \in \Sigma^{\ast}\). From \(q_0\), it moves to all states in \(\ecl\pqty{q_0}\) and only then reads \(\sigma_1\). For each \(x \in \ecl\pqty{q_0}\), it takes the set of states \(\delta\pqty{x, \sigma_1}\), and for each \(y_x \in \delta\pqty{x, \sigma_1}\), it goes to \(\ecl\pqty{y_x}\), and the automaton is now in the set these states, union over all \(y_x\) over all such \(x\). We continue this until the word \(w\) is read, then \(w\) is accepted if any of the end states is an accept state, and rejected otherwise.

\begin{definition}[Extended Transition Function (\(\emptyword\)-NFA)]
    The \term{extended transition function} of an \(\emptyword\)-NFA \(E\) is given by
    \[
        \functiondomain{\hat{\delta}}{Q \times \Sigma^{\ast}}{\powerset\pqty{Q}}
    \]
    defined inductively by
    \begin{itemize}
        \item the empty word \(\hat{\delta} \pqty{q, \epsilon} = \ecl\pqty{q}\),
        \item a one-character word
        \[
            \hat{\delta} \pqty{q, \sigma} = \ecl\pqty{\bigcup_{p \in \ecl\pqty{q}} \delta\pqty{p, \sigma}},
        \]
        and
        \item the inductive step
        \[
            \hat{\delta} \pqty{q, \sigma_1 \ldots \sigma_k} = \ecl\pqty{\bigcup_{p \in \hat{\delta}\pqty{q, \sigma_1 \ldots \sigma_{k - 1}}} \delta\pqty{p, \sigma_k}}.
        \]
    \end{itemize}
\end{definition}

\subsection{Equivalence of DFAs, NFAs and \(\emptyword\)-NFAs}

Clearly, any DFA ``is'' an NFA by defining the new transition function
\[
    \delta'\pqty{q, \sigma} = \Bqty{\delta\pqty{q, \sigma}},
\]
and any NFA is an \(\emptyword\)-NFA, by extending transition for
\[
    \delta\pqty{q, \emptyword} = \emptyset
\]
such that \(\ecl\pqty{S} = S\) for all \(S\).

Underlying all this is, however, languages. Any regular language is the language of an (\(\emptyword\)-) NFA clearly by above.

\begin{definition}[Language Accepted by NFA/\(\emptyword\)-NFA]
    Let \(N\) be an NFA (or an \(\emptyword\)-NFA). The \term{language accepted by \(N\)}, denoted by \(\langacc\pqty{N}\), is the set of all words accepted by \(N\), i.e.,
    \[
        \langacc\pqty{N} = \set{w \in \Sigma^{\ast}}{\hat{\delta}\pqty{q_0, w} \cap F \neq \emptyset}.
    \]
\end{definition}

In fact, the converse is true as well, and we will show this by the \term{subset construction}. We will show this for NFAs first, and the idea for \(\emptyword\)-NFAs naturally follows.

Suppose now we have a NFA \(N\). We will produce a DFA \(D\) on the same alphabet \(\Sigma\), such that \(\langacc\pqty{N} = \langacc\pqty{D}\).

\begin{definition}[Subset Construction]
    \label{def:subset-construction}%
    For the NFA \(N = \pqty{Q, \Sigma, \delta, q_0, F}\), the \term{subset construction} for \(N\) produces the DFA \(D = \pqty{\powerset\pqty{Q}, \Sigma, \Delta, \Bqty{q_0}, \mathcal{F}}\), where the accept states are given by the set of subsets of states that contain at least one accept state of \(N\):
    \[
        \mathcal{F} = \set{S \subseteq Q}{S \cap F \neq \emptyset},
    \]
    and for \(S \subseteq Q, \sigma \in \Sigma\), we define
    \[
        \Delta\pqty{S, \sigma} = \bigcup_{q \in S} \delta\pqty{q, \sigma} \in \powerset\pqty{Q}.
    \]

    (Intuitively, \(\Delta\) gives the subset of states reached from \(S\) by following an edge labelled \(\sigma\).)
\end{definition}

\begin{theorem}
    \label{thm:subset-construction-equivalence}%
    For the NFA \(N\) and DFA \(D\) obtained from the subset construction, and let \(\hat{\delta}\) and \(\hat{\Delta}\) be the respective extended transition functions. Then for all \(w \in \Sigma^{\ast}\), we have
    \[
        \hat{\Delta}\pqty{\Bqty{q_0}, w} = \hat{\delta}\pqty{q_0, w}.
    \]
\end{theorem}

\begin{proof}
    We prove by induction on the length \(k = \abs{w}\). Base case \(k = 0\) is \(w = \emptyword\), then
    \[
        \hat{\Delta}\pqty{\Bqty{q_0}, \emptyword} = \Bqty{q_0} = \hat{\delta}\pqty{q_0, \emptyword}.
    \]

    Now, suppose that
    \[
        \hat{\Delta} \pqty{\Bqty{q_0}, \sigma_1 \ldots \sigma_{k - 1}} = \hat{\delta}\pqty{q_0, \sigma_1 \ldots \sigma_{k - 1}}.
    \]

    Then, for \(w = \sigma_1 \ldots \sigma_{k - 1} \sigma_k\), we have
    \begin{align*}
        \hat{\delta} \pqty{q_0, w} &= \bigcup_{p \in \hat{\delta}\pqty{q_0, \sigma_1 \ldots \sigma_{k - 1}}} \delta\pqty{p, \sigma_k} && \text{by \autoref{def:extended-transition-function-nfa}}\\
        &= \Delta \pqty{\hat{\delta} \pqty{q_0, \sigma_1 \ldots \sigma_{k - 1}}, \sigma_k} && \text{by \autoref{def:subset-construction}}\\
        &= \Delta \pqty{\hat{\Delta} \pqty{\Bqty{q_0}, \sigma_1 \ldots \sigma_{k - 1}}, \sigma_k} && \text{by Induction Hypothesis} \\
        &= \hat{\Delta} \pqty{\Bqty{q_0}, w}. && \text{by \autoref{def:extended-transition-function-dfa}}
    \end{align*}
\end{proof}

\begin{corollary}[Equivalence of DFAs and NFAs]
    A language \(L\) is accepted by a DFA if and only if it is accepted by an NFA.
\end{corollary}

\begin{proof}
    If \(L\) is accepted by a DFA, then it is accepted by an NFA, since any DFA is an NFA.
    
    Conversely, if \(L\) is accepted by an NFA \(N\), then we can construct a DFA \(D\) from the subset construction. Then \(w\) is accepted by \(N\), if and only if
    \[
        \hat{\delta} \pqty{q_0, w} \cap F \neq \emptyset,
    \]
    which by \autoref{thm:subset-construction-equivalence} is if and only if
    \[
        \hat{\Delta} \pqty{\Bqty{q_0}, w} \in \mathcal{F},
    \]
    which is if and only if \(w\) is accepted by \(D\).
\end{proof}